\documentclass[aspectratio=169]{beamer} % Não é obrigatório deixar 16:9

% Desnecessários em compiladores que aceitam UTF-8 por padrão, manter é boa prática
% caso planeje disponibilizar o código
\usepackage[utf8]{inputenc}
\usepackage[brazil]{babel}

% Funções adicionais envolvendo notação matemática
\usepackage{amsmath, amsfonts, amssymb}

% Belezas visuais e funções adicionais de edição
\usepackage{csquotes}
\usepackage{float}
\usepackage{makecell}
\usepackage{dirtytalk} % O nome desse pacote é genial
\usepackage{booktabs}
\usepackage{fancyhdr}
\usepackage{csquotes}

\usepackage{pifont}

\usepackage[document]{ragged2e} % Só pra justificar, não precisa usar isso
\usepackage{xcolor} % serve pra colorir texto
\usepackage{fontspec} % configurações de fonte
\usepackage{multicol} % funções simples de coluna
\usepackage{hyperref} % links externos e internos

% Só para os links ficarem com aparência de link
\hypersetup{
    colorlinks=true,
    linkcolor=blue,
    urlcolor=blue,
    citecolor=black
}


% Referencial bibliográfico em ABNT sem esforço
\usepackage[backend=biber,style=abnt]{biblatex} % abntex2cite precisa de "root"

\setbeamersize{text margin left=0.8cm, text margin right=0.8cm}

% Endereço da pasta de imagens e arquivos adicionais
\graphicspath{{./Extras/}}

% Configuração de fonte
\setmainfont{Times New Roman}

% Endereço do arquivo de referências
\addbibresource{referencias.bib}

% Vários temas ficam disponíveis, uma busca rápida no google mostra opções, a
% escolha fica a seu critério.
\usetheme{Boadilla}


% Cabeçalho interno
\title{Título da apresentação}

\subtitle{subtítulo}
\author{Autoria}
\date{\today}



% A partir daqui começamos o documento propriamente dito
\begin{document}

% Abaixo o exemplo de geração de slide título, este usa o cabeçalho fornecido
\frame{\titlepage}


\begin{frame}

    \frametitle{Introdução}

    \begin{center}
        ``DON'T PANIC"
       \textit{\cite{adams}}
    \end{center}

    \begin{block}{\centering O que é \LaTeX?}
        \centering
        Trata-se de um sistema de preparação de documentos científico/acadêmicos.

    \end{block}

    Para quem nunca teve contato com linguagens de programação ou correlatos, a
    primeira experiência com a linguagem pode ser intimidadora. Confie na gente, é
    bem mais simples do que a sintaxe pomposa faz parecer. Links úteis pra conferir
    (só clicar):
    \vspace{6pt}

    \begin{itemize}

        \item \href{https://learnxinyminutes.com/pt-br/latex/}{Guia de primeiros
            passos em português (recomendação pessoal)}

        \item \href{https://www.latex-project.org/help/documentation/\#general-documentation}{Documentação
            oficial}

        \item \href{https://www.overleaf.com/learn/latex/Learn\_LaTeX\_in\_30\_minutes}{Guia
            do Overleaf}

        \item \href{https://typst.app/docs/reference/symbols/sym/}{Símbolos
            matemáticos}

    \end{itemize}


\end{frame}



\begin{frame}
    \frametitle{Funcionamento básico}

    A compilar, o \LaTeX segue uma série de parâmetros pra otimizar e uniformizar
    textos acadêmicos. Quando se entende a lógica base tudo fica bem simples.

    No caso de aprestações, estruturas extras e convenções
    adicionais de sintaxe são usadas. Tentamos incluir a maioria delas nesse slide
    exemplo.

    Abaixo exemplos de bullet:
    \vspace{1em} % insere espaçamento horizontal


    % A primeira vista isso pode parecer confuso ou prolixo, de verdade é mais
    % simples do que parece. A maior parte da sintaxe é preenchida por auto complete
    % pelos editores de código mais populares a disposição.

    % Minipages e disposição de colunas serão abordadas no capítulo seguinte

    \begin{multicols*}{2}

        \begin{minipage}{\columnwidth}
            \begin{itemize}

                \item Aqui consta um bullet
                \item Aqui consta outro bullet
                \item Advinha? Mais um bullet

            \end{itemize}

        \end{minipage}


        \columnbreak % quebra colunas

        \begin{minipage}{\columnwidth}

            \begin{enumerate}

                % Funciona dos dois jeitos abaixo, delimitação explícita ou implícita
                \item{Essa é outra maneira de criar bullets}
                \item Bullet fundamental
                \item Imagine algo importante

            \end{enumerate}
        \end{minipage}


    \end{multicols*}


\end{frame}


\begin{frame}
    \frametitle{Colunas e disposição de elementos}

    \begin{multicols*}{2} % Define o número de colunas utilizado no campo

        \begin{minipage}{\columnwidth}
            \hfill % Comando útil pra lidar com controle de espaçamentos horizontais

            Aqui criamos um ambiente segmentado em duas colunas.
            O asterisco no ambiente e as \emph{minipages} server pra que o \LaTeX
            não tome liberdades na disposição do conteúdo.

            \begin{block}{\centering Observação}
                Leia o código junto com o PDF, nele há comentários e exemplos que
                facilitam a compreensão.

            \end{block}

            \hfill

        \end{minipage}

        \columnbreak

        \begin{minipage}{\columnwidth}

            \begin{figure}[H]

                \centering

                % O href foi somente pra que o qr code fosse clicável e abrisse o
                % link. É bom evitar usar isso porque redirecionamento de links é um
                % comportadamente que pode ser interpretado como potencialmente
                % malicioso por alguns antivirus, porque abre margem pra alguém por
                % links maliciosos. Note

                \href{https://www.overleaf.com/learn/latex/Inserting\_Images\%23Positioning}{
                \includegraphics[width=0.3\textwidth]{over_images.png}
            }

            \caption{Documentação de exemplos com imagens}
            \label{fig:hehe}

            \end{figure}

            \vspace{-1em}

            % usasse "em" pra se referir ao tamanho padrão de uma linha, vspace
            % adiciona ou remove espaços verticalmente. Também pode-se usar pt, que
            % refere-se a tamanho de fonte e cm pra referir-se ao tamanho diretamente
            % em centímetros

            Ao citar objetos usa-se um \emph{label}, fazer isso
            garante que quaisquer mudanças sejam reconhecidas na hora
            de compilar o documento. Ordenamentos ficam corrigidos assim, exemplo:
            figura \ref{fig:hehe}.


        \end{minipage}

    \end{multicols*}


\end{frame}



\begin{frame}
    \frametitle{Equações e notação matemática}

    Exemplo integrado de equação:
    ``Uma vez o físico da língua pra fora disse: $E = MC^2$"

    \vspace{1em}
    Exemplo solto de equação:

    % Dá pra fazer isso com símbolos de cifrão também, fica a critério
    \[
            Z[J]=
            \int \mathcal{D}\phi\,
            \exp\left\{
                i\int_{\mathbb{R}^{1,3}}
                \left[
                    \frac{1}{2}\partial_\mu\phi\,\partial^\mu\phi
                    -\frac{m^2}{2}\phi^2
                    -\frac{\lambda}{4!}\phi^4
                    +J(x)\phi(x)
                    \right]d^4x
                \right\}
        \]

        \vspace{1em}
        Exemplo de ambiente de equação:


        \begin{equation}
            \begin{aligned}
                \underbrace{
                    \frac{\partial \mathbf{u}}{\partial t}
                    +(\mathbf{u}\cdot\nabla)\mathbf{u}
                    =-\frac{1}{\rho}\nabla p
                    +\nu\nabla^2\mathbf{u}
                    +\mathbf{f},
                }_{\text{Navier-Stokes}}
                \\[6pt]
                \nabla\cdot\mathbf{u}&=0.
            \end{aligned}
            \label{eq:navier}
        \end{equation}




\end{frame}



\begin{frame}
    \frametitle{Tabelas}
    % O conteúdo a seguir foi feito pelo stargazer (R)

    % Aqui consta a função padrão para separar colunas, citamos a outra por que
    % interage melhor com alguns tipos de edição
    \begin{columns}

        \column{0.6\textwidth}

        \begin{table}[!htbp] \centering
            \caption{Reressão meramente ilustrativa}
            \label{}
            \begin{tabular}{@{\extracolsep{5pt}}lc}
                \\[-1.8ex]\hline
                & \multicolumn{1}{c}{\textit{Dependent variable:}} \\
                \cline{2-2}
                \\[-1.8ex] & agro \\
                \hline \\[-1.8ex]
                Agro & 0,176$^{***}$ \\
                & (0,002) \\
                & \\
                Pib & $-$0,0001$^{***}$ \\
                & (0,00003) \\
                & \\
                \hline \\[-1.8ex]
                Observations & 103.189 \\
                R$^{2}$ & 0,114 \\
                \hline \\[-1.8ex]
                \textit{Note:}  & \multicolumn{1}{r}{$^{*}$p$<$0,1; $^{**}$p$<$0,05; $^{***}$p$<$0,01} \\
            \end{tabular}
        \end{table}

        \column{0.4\textwidth}

        Exemplo de colunas assimétricas e uma tabela de regressão gerada no stargazer
        do R.

        {
            \Huge
            \vspace{-1.2em}
            $$\leftarrow$$
        }

        Tabelas podem gerar alguma dificuldade inicial. Linguagens de programação tem
        ferramentas que já as formatam pra LaTeX.

        \begin{itemize}
            \item \href{https://www.tablesgenerator.com/latex\_tables}{Ferramenta
                online pra converter tabelas CSV ou excel para LaTeX}

            \item \href{https://www.overleaf.com/learn/latex/Tables}{Documentação do
                Overleaf}

        \end{itemize}

        % citação fantasma
        \nocite{latexproject}
        \nocite{latex_in_y}
        \nocite{overleaf_learn_latex}


    \end{columns}


\end{frame}

\begin{frame}

    \begin{center}
        \Huge
        Valeu, pessoal!
    \end{center}

\end{frame}


\begin{frame}
    \frametitle{Referências}
    \printbibliography

\end{frame}


\end{document}

